\documentclass[10pt,a4paper]{amsart}
\usepackage[margin=19mm]{geometry}
\usepackage{amsmath,amssymb,mathtools}
\usepackage[scr=rsfs,cal=euler]{mathalfa}
\usepackage{xfrac}
\usepackage[unicode,hidelinks,bookmarks=false]{hyperref}
\newcommand{\ie}{i.e.\ }
\newcommand{\cf}{cf.\ }
\newcommand{\resp}{respectively\ }
\newcommand{\Subsection}{\subsection}
\providecommand{\nobreakdash}{\nobreak\hbox{-}}
\setlength{\emergencystretch}{2em}
\multlinegap0pt
\usepackage{amssymb}
\usepackage{float}
\usepackage[shortlabels]{enumitem}
\newtheorem{lemma}{Lemma}
\newtheorem{theorem}{Theorem}
\title{Complexity of Linear Subsequences of Fibonacci-Automatic Sequences}
\author{Delaram Moradi, Narad Rampersad, Jeffrey Shallit}
\date{Source: arXiv:2603.21645v1 (CC BY 4.0)}
\begin{document}
\maketitle

\noindent\emph{Source and licence.} Complexity of Linear Subsequences of Fibonacci-Automatic Sequences. Version arXiv:2603.21645v1 (CC BY 4.0), distributed by arXiv under the Creative Commons Attribution 4.0 International licence, http://creativecommons.org/licenses/by/4.0/, as recorded in the arXiv licence metadata for this exact version and shown on the abstract page https://arxiv.org/abs/2603.21645v1. Full licence notice: \texttt{licenses/arxiv-2603.21645-CC-BY-4.0.txt}.\par\smallskip

\noindent\textit{Editorial scope.} Counted unit: the article's theorem theorem:lin-subseq-sc (source0.tex lines 729--734). The article's complete original proof of it is delivered with it (source0.tex lines 793--831, one \texttt{proof} environment of 39 lines, which the article places away from the statement). No mathematics is written, repaired or re-derived: the statement and the proof are the article's own lines, in the article's own notation, and no retained line is substituted or re-flowed.  Retained uncounted as the source-local context the proof needs: lines 189--194; lines 297--301; lines 773--780.  Nothing was omitted from any retained span, so the package carries no omission record.  No retained line contains a citation, so the wrapper carries no bibliography and no bibliography entry.  No retained line contains a figure environment, an image-inclusion command or drawing code, so no image is delivered and the package declares no asset dependency.  Editorial formatting only, in addition: an AMS \texttt{amsart} preamble that restates the article's own declarations and macros used by the retained lines, namely the environments \texttt{lemma}, \texttt{theorem} on separate counters, together with the standard packages those lines need.  The retained environments print in this document as Lemma 1, Theorem 1 in order of appearance, whereas the article prints the same items with the numbering of its own full document.  The complete native source of the article as distributed in the arXiv e-print for this exact version is bundled unchanged as \texttt{sources/196964/source0.tex}.
\begin{lemma}
We have $\rho_{\bf x}(n) = O(n)$.
\label{lemma:subword-sc}
Let ${\bf x}$ be  Fibonacci-automatic sequence generated by an $m$-state DFAO with msd-first input.
Then $\rho_{\bf x} (n) = O(n m^2)$ for all $n \geq 1$. 
\end{lemma}

\begin{lemma}
Suppose $x$ is a word, $a \in \{0,1\}$ and $xa$ is a valid Fibonacci representation.
Then $[x] = \lfloor ([xa]+2)/\alpha \rfloor - 1$.
\label{lemma:[x]-floor}
\end{lemma}

\begin{lemma}
For each state $p$ in $M$
there exist $a \in \Sigma_2$ and integers $S, T$ such 
that $[S, S+T] \subseteq [-n,2n-1]$ and every state of $p$ is of the form
$[a, b_i, d_i, d'_i, s_i]$ 
with $0 \leq i \leq T$ and $d_i = S+i$.
\label{lemma:lin-subseq-consec}
\end{lemma}

\begin{theorem}
Let $n \geq 1$ and $0 \leq c < n$, and let $(h(i))_{i \geq 0}$ be a Fibonacci-automatic sequence generated by a DFAO of $m$ states.  There is a DFAO of
$O(m^2n^4)$ states recognizing the linear subsequence
$(h(ni+c))_{i \geq 0}$.
\label{theorem:lin-subseq-sc}
\end{theorem}

\begin{proof}[Proof of Theorem \ref{theorem:lin-subseq-sc}]
Consider an input $xa$ to $M$.
By Lemma~\ref{lemma:lin-subseq-consec} we know that
on this input $M$ reaches a set of
states corresponding to inputs $xa\times y_i b_i$ in $M_u$, where $d_i = [y_i b_i]-n[xa]=  S+i$  for
$0 \leq i \leq T$, some word $y_i$ and letter $b_i$.  Let $[y_0 b_0] = Y$.  Then $[y_i b_i] = Y+i$
for $0 \leq i \leq T$.  Noting that
$d'_i = [y_i] - n[x]$, 
define $e_i = d'_{i+1} - d'_i$ for $0 \leq i < T$. 
Let $g$ be the function that
maps $t$ to the integer obtained by dropping the last
bit in the Fibonacci representation of $t$; Lemma~\ref{lemma:[x]-floor}
shows that $g(t) = \lfloor (t+2)/\alpha \rfloor - 1$.
Then $d'_i = g(Y+i) -n[x]$.
The triples $(b_i, e_i, s_i )$ are given by $({\bf f}[Y+i], g(Y+i+1)-g(Y+i), h'(Y+i))$
for $0 \leq i \leq T$, and this formula completely specifies {\it all\/} the $b_i$ and $d'_i$ and $s_i$,
once $d'_0$ is chosen.  Namely,
$d'_i = d'_0 + \sum_{0 \leq j < i} e_j$.

On input $x$ the automaton generating $(g(j+1)-g(j))_{j \geq 0}$ only needs to keep track of the last two letters read; if any of them is $1$ then $g([x]) = 1$, otherwise $g([x]) =0$. 
Now think of the triples $(b_i, e_i, s_i)$  as a subsequence of the infinite sequence of triples 
${\bf x} := ({\bf f}[j], g(j+1)-g(j), h'(j))_{j\geq 0}$, which 
is an automatic sequence because $({\bf f}[j])_{j \geq 0}$, $(g(j+1)-g(j))_{j \geq 0}$, and the interior sequence $(h'(j))_{j\geq 0}$ are automatic sequences.
Furthermore, the sequence ${\bf x}$ is generated by a DFAO of $O(m)$ states.
The finite sequence
$({\bf f}[Y+i], g(Y+i+1)-g(Y+i), h'(Y+i))_{0 \leq i \leq T}$ is then completely specified by saying which factor of length $T+1$ of $\bf x$ it is, and by
Lemma~\ref{lemma:subword-sc}, there are $O(m^2n)$ of them.

Summarizing, there are
\begin{itemize}
\item[(a)] $O(n)$ choices for $d_0 \in [-n, 2n-1]$,
\item[(b)] $O(n)$ choices for $d'_0 \in [-n, 2n-1]$,
\item[(c)] $O(n)$ choices for $T$,
\item[(d)] $O(m^2n)$ choices for which particular subword of the automatic sequence $\bf x$ gives $(b_i, d'_{i+1}-d'_i, s_i)$ for $0 \leq i \leq T$,
\item[(e)] $O(1)$ choices for $a$.
\end{itemize}

Therefore there are $O(m^2n^4)$ possible states.
\end{proof}

\end{document}
